\chapter{MAC-IO (Proposed): Self-Supervised Cross-Platform Adaptation for Inertial Odometry} \label{chapter:macio}
\fancyhead[LO]{\makebox[\textwidth][r]{Chapter 9. MAC-IO (Proposed)}}

\section{Introduction}

This chapter proposes MAC-IO, a self-supervised adaptation framework for learned inertial odometry.
The starting point is the completed AirIO model in Chapter~\ref{chapter:airio}, which provides strong inertial odometry performance but can overfit to its training-domain motion and sensor statistics.
When deployed on new platforms, mounting conditions, and unseen trajectories, accuracy can degrade due to distribution shift.

MAC-IO addresses this problem by leveraging the metrics-aware inertial uncertainty and preintegration formulation established in AirIMU (Chapter~\ref{chapter:airimu}) as additional supervision signals.
The key idea is to adapt the AirIO velocity network on new data without requiring external ground truth.
To achieve this goal, MAC-IO proposes two self-supervision pathways that are complementary in efficiency and supervision strength.

The motivation is that collecting ground-truth motion labels for each new platform is expensive and often impractical.
This limitation is especially severe for aerial, off-road, and safety-critical deployments where data collection conditions are difficult to control.
As a result, many learning-based inertial odometry systems are accurate in-domain but fragile out-of-domain.
MAC-IO targets this bottleneck by converting estimator structure itself into supervision, so adaptation can occur directly on unlabeled target-domain data.

The chapter is designed as an adaptation layer on top of prior completed work.
AirIMU provides metrics-aware inertial uncertainty and preintegration consistency.
AirIO provides a strong inertial velocity estimator and EKF fusion loop.
MAC-IO combines these ingredients to create a practical transfer mechanism that favors fast large-scale adaptation when compute is limited and stronger iterative refinement when filter feedback is available.

\section{Proposed Method}
\label{sec:macio:method}

\subsection{Problem Setup}

Let a pretrained inertial odometry network predict velocity $\hat{\mathbf{v}}_t$ from IMU windows.
Given target-domain IMU data without labels, MAC-IO adapts network parameters $\theta$ so that inertial predictions remain consistent with physically grounded estimation signals.
The adaptation objective combines two self-supervised losses:
\begin{equation}
\label{eq:macio:total}
\mathcal{L}_{\text{MAC-IO}} =
\lambda_{\Delta v}\,\mathcal{L}_{\Delta v}
+
\lambda_{\text{EKF}}\,\mathcal{L}_{\text{EKF}}.
\end{equation}

\subsection{Method 1: IMU Preintegration Delta-V Supervision}

The first method uses IMU preintegration as a fast self-supervision signal.
From raw inertial measurements, we compute the preintegrated velocity increment $\Delta\mathbf{v}^{\text{imu}}_{t\rightarrow t+1}$.
From network outputs, we form the predicted increment $\Delta\hat{\mathbf{v}}_{t\rightarrow t+1} = \hat{\mathbf{v}}_{t+1} - \hat{\mathbf{v}}_{t}$.
The adaptation loss is:
\begin{equation}
\label{eq:macio:deltav}
\mathcal{L}_{\Delta v} =
\sum_{t}
\left\|
\Delta\hat{\mathbf{v}}_{t\rightarrow t+1}
-
\Delta\mathbf{v}^{\text{imu}}_{t\rightarrow t+1}
\right\|_{\Sigma^{\Delta v}_{t}}^{2},
\end{equation}
where $\Sigma^{\Delta v}_{t}$ is derived from metrics-aware inertial uncertainty.

This method is computationally light, easy to scale, and suitable for large target-domain datasets.
Its limitation is that it supervises velocity \textit{changes} rather than absolute velocity, so bias in global velocity can persist.

\subsection{Method 2: EKF-Filtered Velocity Supervision}

The second method uses the filtered state from an uncertainty-aware EKF as a stronger supervisory target.
After fusing AirIO network outputs with IMU dynamics in the EKF, we obtain filtered states $(\mathbf{p}^{f}_t, \mathbf{v}^{f}_t, \mathbf{R}^{f}_t)$.
MAC-IO then supervises network velocity outputs directly with filtered velocity:
\begin{equation}
\label{eq:macio:ekf}
\mathcal{L}_{\text{EKF}} =
\sum_{t}
\left\|
\hat{\mathbf{v}}_{t}
-
\mathbf{v}^{f}_{t}
\right\|_{\Sigma^{f}_{t}}^{2},
\end{equation}
where $\Sigma^{f}_{t}$ is the EKF state covariance used as confidence weighting.

This method follows a bilevel-style adaptation intuition: the network proposal is fused into the filter, and the improved filtered state is fed back as supervision for iterative refinement.
Compared with Method 1, this pathway can supervise absolute velocity and typically yields stronger adaptation, at the cost of higher computation and tighter coupling with the filter loop.

\subsection{Expected Adaptation Behavior}

Together, the two methods provide a practical adaptation hierarchy:
\begin{itemize}
    \item \textbf{Method 1} for fast large-scale adaptation with lightweight supervision.
    \item \textbf{Method 2} for stronger iterative refinement using filter-improved state targets.
\end{itemize}
The expected outcome is improved inertial odometry performance on unseen platforms and trajectories without manual retuning or external ground-truth labels, directly supporting the thesis objective of minimal additional human engineering.